% introduction.tex
% Orientation for Paper IV: the lattice as a birefringent medium with a
% geometry-fixed optical axis; the two channels (kinematic + gauge); the
% companion structure with Paper VIII; scope. The audit table is included in
% the Scope subsection so the reader meets it before any status-dependent claim.

\label{sec:introduction}

A crystal is birefringent when its atomic lattice singles out an axis, so that the
speed of light depends on direction and polarization. The A=1 discrete causal
lattice of Paper~I~\cite{menendez2026geometryfirst} is birefringent for the same
reason, but the ``crystal'' is spacetime itself. Its bipartite hop set---the three
RGB vectors $\mathbf{V}_1=(1,1,1)$, $\mathbf{V}_2=(1,-1,-1)$,
$\mathbf{V}_3=(-1,1,-1)$ and their CMY negatives---uses three of the four cube
body-diagonals and leaves the fourth, $\hat{\mathbf{n}}_* = (1,1,-1)/\sqrt3$, out.
That single omission breaks the cubic point group $O_h$ to the trigonal $D_{3d}$
about $\hat{\mathbf{n}}_*$ and fixes an optical axis geometrically, with no new
postulate. The optical-axis birefringence is one of the series' standing
predictions; this paper works it to closed form.

The effect appears in two sectors, and they share the \emph{same} axis. The
\emph{kinematic} channel (Sec.~\ref{sec:kinematic_channel}) follows from Paper~I's
single-tick spinor propagator: the entire wavevector dependence passes through one
scalar structure factor whose modulus is flat along $\hat{\mathbf{n}}_*$ and curved
perpendicular to it, giving a closed-form ordinary/extraordinary speed split. This
is the falsifiable, direction-fixed dispersion prediction the paper contributes, and
its observability---that the anisotropy survives the per-tick A=1 renormalization
into observable amplitudes---is established numerically. The \emph{gauge} channel
(Sec.~\ref{sec:gauge_sector}) is the response of the vacuum: integrating the matter
tokens out of the dynamics in a background electromagnetic field induces an
anisotropic Maxwell action whose permittivity $\boldsymbol{\varepsilon}$ and
inverse permeability $\boldsymbol{\mu}^{-1}$ are again uniaxial about
$\hat{\mathbf{n}}_*$. Here the outcome is the opposite of a new signal: the two
tensors obey the adjugate relation $\boldsymbol{\mu}^{-1} =
\operatorname{adj}(\boldsymbol{\varepsilon})$, which forces the two photon
polarizations to share a dispersion, so the gauge-sector vacuum \emph{birefringence
cancels}. What remains is a common-mode speed anisotropy---a direction-dependent
vacuum speed affecting both polarizations equally---whose observational status is
left open. That the two sectors inherit one axis from one geometry is the paper's
multi-channel concordance.

The gauge-sector story is developed analytically in the companion
Paper~VIII~\cite{menendez2026electricblock}, which derives the two blocks, proves the
adjugate cancellation theorem, and settles the electric sector that neither Paper~I
nor Paper~II supplied. The present paper is the independent, from-engine confirmation:
both blocks are read off the running lattice's hop holonomies, the cancellation is
verified numerically on those engine-extracted tensors, and the two papers are
intended as a back-to-back pair.

\subsection{Scope}
\label{subsec:scope}

In scope, and derived here: the closed-form kinematic speed split
(Sec.~\ref{sec:kinematic_channel}, \texttt{PASS}); the from-engine extraction of the
gauge-sector blocks and the null birefringence verdict
(Sec.~\ref{sec:gauge_sector}, \texttt{PASS}); and the structural axis-sharing of the
two channels. Explicitly deferred: the isotropic overall scale of the induced action
(the coupling $1/g^2$ Paper~I leaves open), so the gauge sector fixes anisotropic
\emph{structure} but not the absolute vacuum speed; the observational status of the
common-mode speed anisotropy (the single-domain-versus-$O_h$-restoration question,
\texttt{PART}); and the single-object \emph{dynamical} ($\operatorname{Tr}\ln T$)
extraction of the induced-action tensor, which is a stated open item
(Remark~\ref{rem:trlnt_open}). The quantitative falsifiability of the kinematic
prediction---its magnitude against current astrophysical bounds and its mapping onto
the Standard-Model Extension---is analyzed in the accompanying notes rather than
claimed in the prose. The audit table (Table~\ref{tab:audit}) is the canonical record
of which claims are \texttt{PASS}, \texttt{PART}, \texttt{STUB}, or \texttt{FAIL}.

% ============================================================
% audit_table.tex
%
% A "system audit" table mapping the paper's claims to the evidence
% backing each one.  This is *the* canonical source of truth for which
% claims are PASS / PART / STUB / FAIL --- prose elsewhere in the paper
% should not contradict it.
%
% Status semantics (paste into the abstract or front-matter as well):
%   PASS  -- derivation complete or experiment confirms the claim to
%            stated precision.
%   PART  -- mechanism demonstrated, full quantitative match pending;
%            specific gaps documented in the evidence column.
%   STUB  -- analytical result with no numerical confirmation yet,
%            or planned experiment not yet run.
%   FAIL  -- tested and disconfirmed.
%
% Use \texttt{} for code/experiment names and file paths inside cells.
% Long cells are fine -- the column widths are tuned for prose.
% ============================================================

\label{tab:audit}

{\scriptsize
\setlength{\tabcolsep}{4pt}
\renewcommand{\arraystretch}{0.95}
\begin{longtable}{@{}p{2.6cm}p{3.8cm}p{3.2cm}p{4.4cm}p{1.0cm}@{}}
\caption{\small System audit: paper claims mapped to supporting
derivations and experiments.  Status semantics in the front-matter
(\texttt{PASS} / \texttt{PART} / \texttt{STUB} / \texttt{FAIL}).}
\label{tab:audit_caption} \\
\hline
\textbf{Claim} & \textbf{Mechanism} & \textbf{Continuum / formal limit} & \textbf{Evidence} & \textbf{Status} \\
\hline
\endfirsthead

\multicolumn{5}{l}{\textit{(Table~\ref{tab:audit}, continued from previous page)}} \\
\hline
\textbf{Claim} & \textbf{Mechanism} & \textbf{Continuum / formal limit} & \textbf{Evidence} & \textbf{Status} \\
\hline
\endhead

\hline
\multicolumn{5}{r}{\textit{(continued on next page)}} \\
\endfoot

\hline
\endlastfoot

% ── Kinematic channel ───────────────────────────────────────────────
Kinematic birefringence along $(1,1,-1)$
    & Closed-form uniaxial split from the spinor propagator: $M_\text{eff} = \tfrac43(\mathbb{I}-\hat{\mathbf{n}}_*\hat{\mathbf{n}}_*^T)$, so $|H_\text{RGB}(\mathbf{k})|^2 = 1-\tfrac43|\mathbf{k}_\perp|^2$. Group-velocity split $\Delta v = \tfrac23\sin\omega\,|\mathbf{k}|$ (angular law $\sin\alpha$, massive sector) and $\omega$-independent dispersion-curvature split $4/3$ vs $0$ (photon + massive)
    & Isotropic $c$ under $O_h$ averaging / $a\to 0$
    & \emph{Analytic:} eigenphase dispersion $\tan\theta = \tan(\omega/2)/|H_\text{RGB}|$ and group velocity in closed form; $|H_\text{RGB}|^2 = 1$ exactly along $(1,1,-1)$ at all orders. \texttt{exp\_02\_closed\_form\_split\_check}: structure-factor and dispersion identities to $<10^{-12}$, speed coefficients to $<10^{-5}$. \emph{Observability:} \texttt{exp\_01\_optical\_axis\_isotropy} ($129^3$, both engines) -- isotropic pulse flattens perpendicular to $(1,1,-1)$ (ratio $\to 0.01$), flat-axis alignment $=1.0000$, $\omega$-independent, survives exact A=1 renormalization.
    & \texttt{PASS} \\

% ── Gauge sector ────────────────────────────────────────────────────
Gauge-sector birefringence ($\mathbf{Q}$ eigenvalues $\{4,4,16\}$): uniaxial structure derived; observability \emph{open}
    & Induced photon action $\tfrac{ca^4}{2}\mathbf{F}^T\mathbf{Q}\mathbf{F}$; special direction Hodge-duals to the same axis $(1,1,-1)$. The bipartite RGB/CMY structure breaks $O_h$ to a uniaxial subgroup about $(1,1,-1)$ ($\mathbf{Q}$ fixed by $12/48$, RGB by $6/48$), so the action is genuinely uniaxially anisotropic
    & $O_h$-averaged Maxwell is the \emph{orientation-averaged} form, not what a fixed lattice's observables see; quantitative $\mathbf{Q}$ role is the coupling ratio $g_3^2/g_2^2 = 3/2$ (Paper~II)
    & \emph{Structure derived} (Paper~I App.~B; Paper~II). \emph{Observability OPEN:} no symmetry forces isotropy (the $O_h$-average is an external imposition the uniaxial dynamics do not justify); by the \texttt{exp\_01} analogy the anisotropy may survive in observables, where read literally it is a dimension-4, $O(1)$ effect excluded by $\sim$30--37 orders. The one possible escape -- electric/magnetic cancellation in the full dispersion (Paper~I did the magnetic sector only) -- is uncomputed. \texttt{exp\_03} (full E+B Peierls dynamics) is decisive. See \texttt{notes/gauge\_sector\_structural\_conclusion.md}.
    & \texttt{PART} \\

% ── Concordance ─────────────────────────────────────────────────────
Multi-channel concordance (P9)
    & Kinematic and gauge sectors inherit the \emph{same} optical axis $(1,1,-1)$ from the common bipartite geometry ($M_\text{eff}$ and $\mathbf{Q}$ special directions coincide)
    & N/A (cross-sector structural identity)
    & Axis-sharing is a \emph{structural} identity (derived): $M_\text{eff}$ and $\mathbf{Q}$ both inherit $(1,1,-1)$ from the same bipartite geometry. The kinematic dispersion is the one frontier-testable observational channel; the gauge sector's observable status is itself \emph{open} (see its row -- it may be unobservable, suppressed, or an excluded $O(1)$ effect, pending \texttt{exp\_03}). P9's strength as a multi-witness empirical claim is therefore hostage to that outcome.
    & \texttt{PART} \\

\hline
\end{longtable}
}

